\section{Μεθοδολογία}

Η μεθοδολογία της εργασίας συνδύασε θεωρητική μελέτη, αρχιτεκτονική σχεδίαση και επαναληπτική υλοποίηση. Αρχικά πραγματοποιήθηκε 
μελέτη των προδιαγραφών \ENG{SCORM}, του ιστορικού πλαισίου των προτύπων \ENG{E-Learning} και των βασικών χαρακτηριστικών υπαρχουσών λύσεων, 
εμπορικών και ανοιχτού κώδικα. Η φάση αυτή βοήθησε στον προσδιορισμό των λειτουργικών απαιτήσεων του συστήματος και των βασικών αρχιτεκτονικών
προβλημάτων που εμφανίζονται σε πραγματικές υλοποιήσεις.

Η ανάπτυξη του συστήματος ακολούθησε επαναληπτική προσέγγιση εμπνευσμένη από πρακτικές \ENG{Agile} (και συγκεκριμένα του \ENG{SCRUM Framework}).
Αρχικά υλοποιήθηκε το ίδιο \ENG{\texttt{engine}} που αποτελεί τον πυρήνα του συστήματος, ο οποίος περιλαμβάνει τα βασικά μοντέλα δεδομένων και τους μηχανισμούς διαχείρισης.
Στη συνέχεια αναπτύχθηκε ο \ENG{\texttt{player}}, ο οποίος αναλαμβάνει την εκτέλεση του περιεχομένου και την επικοινωνία με τον \ENG{engine}.
Τέλος δημιουργήθηκε το \ENG{\texttt{PHP SDK}} καθώς και ένα ενδεικτικό \ENG{LMS}
(\ENG{\texttt{example-lms-client}}) σε \ENG{Laravel}, με στόχο την επίδειξη της ενοποίησης του συστήματος.

Η συγγραφή της εργασίας εξελίχθηκε παράλληλα με την υλοποίηση ώστε οι αρχιτεκτονικές αποφάσεις να τεκμηριώνονται κατά τη διάρκεια της ανάπτυξης.
Στη διαδικασία υποστηρικτικά χρησιμοποιήθηκαν εργαλεία τεχνητής νοημοσύνης για ιδεασμό, γλωσσική επιμέλεια και βοήθεια σε δευτερεύοντα τμήματα
κώδικα, όπως τα \ENG{OpenAI ChatGPT}, \ENG{OpenAI Codex} και εργαλεία της \ENG{Genspark}. Οι τελικές τεχνικές επιλογές, 
η επιβεβαίωση της ορθότητας και η διαμόρφωση του κειμένου πραγματοποιήθηκαν από τον συγγραφέα.
